\documentclass[a4paper,11pt]{article}
\usepackage[utf8]{inputenc}
\usepackage[french]{babel}
\usepackage[T1]{fontenc}
\usepackage{amsmath,amssymb,amsthm}
\usepackage{geometry}
\geometry{margin=2cm}
\usepackage{enumitem}
\usepackage{hyperref}
\usepackage{booktabs}
\usepackage{xcolor}
\usepackage{listings}
\usepackage{titlesec}
\usepackage{fancyhdr}
\usepackage{lastpage}
\AtBeginDocument{\shorthandoff{;:!?}}

\titleformat{\section}{\large\bfseries\color{blue!60!black}}{}{0pt}{}[\vspace{-0.5ex}\rule{\textwidth}{0.4pt}]
\titleformat{\subsection}{\bfseries\color{blue!40!black}}{}{0pt}{}

\lstset{basicstyle=\ttfamily\small,breaklines=true,frame=single,aboveskip=6pt,belowskip=6pt}

\newtheorem*{important}{\`A retenir}
\pagestyle{fancy}
\fancyhf{}
\fancyhead[L]{\small STA211 -- Compl\'ement Tanagra Stacking Python}
\fancyhead[R]{\small Fiche r\'ecapitulative}
\fancyfoot[C]{\thepage/\pageref{LastPage}}
\renewcommand{\headrulewidth}{0.4pt}

\title{\textbf{Fiche r\'ecapitulative -- STA211}\\[4pt]
\large Compl\'ement \texttt{fr\_Tanagra\_Stacking\_Python} \\[2pt]
\small Stacking -- Agr\'egation par m\'eta-apprentissage avec Python}
\author{Tanagra Data Mining -- R. Rakotomalala}
\date{15 juin 2026}

\begin{document}
\maketitle
\thispagestyle{empty}

\begin{abstract}
Le stacking combine plusieurs classifieurs h\'et\'erog\`enes via un
m\'eta-classifieur qui apprend \`a corriger leurs erreurs. Ce tutoriel
l'impl\'emente en Python avec scikit-learn sur un probl\`eme de
discrimination binaire dans le plan.
\end{abstract}

\vfill
\tableofcontents
\newpage

% ===================================================================
\section{Donn\'ees}
% ===================================================================

Probl\`eme de discrimination binaire dans le plan. 2000 observations,
2 variables pr\'edicatives $X_1, X_2$. Fronti\`ere parabolique~:

\[
\text{Si } 0.1 \times X_2 > X_1^2 \implies \text{pos, sinon neg}
\]

Pr\'eparation~:
\begin{lstlisting}
import numpy as np
import pandas as pd
from sklearn.model_selection import train_test_split

D = pd.read_csv("artificial2d_data2.txt", sep="\t")
y = np.where(D["Y"] == "pos", 1, 0)
X = D[["X1", "X2"]].values

X_train, X_test, y_train, y_test = train_test_split(
    X, y, train_size=1500, test_size=500,
    random_state=100
)
\end{lstlisting}

\noindent 1500 observations en apprentissage, 500 en test.

% ===================================================================
\section{Classifieurs de base}
% ===================================================================

\subsection{R\'egression logistique}

\begin{lstlisting}
from sklearn.linear_model import LogisticRegression
lr = LogisticRegression(max_iter=5000)
lr.fit(X_train, y_train)
print(f"LR: {lr.score(X_test, y_test):.3f}")
\end{lstlisting}

\noindent Accuracy test~: \textbf{0.678} ($\approx$67.8\,\%).

\subsection{K plus proches voisins (KNN)}

\begin{lstlisting}
from sklearn.neighbors import KNeighborsClassifier
knn = KNeighborsClassifier(n_neighbors=5)
knn.fit(X_train, y_train)
print(f"KNN: {knn.score(X_test, y_test):.3f}")
\end{lstlisting}

\noindent Accuracy test~: \textbf{0.918}.

\subsection{Arbre de d\'ecision}

\begin{lstlisting}
from sklearn.tree import DecisionTreeClassifier
tree = DecisionTreeClassifier(max_depth=5, random_state=0)
tree.fit(X_train, y_train)
print(f"Tree: {tree.score(X_test, y_test):.3f}")
\end{lstlisting}

\noindent Accuracy test~: \textbf{0.942}.

\subsection{SVM (noyau RBF)}

\begin{lstlisting}
from sklearn.svm import SVC
svm = SVC(probability=True, random_state=0)
svm.fit(X_train, y_train)
print(f"SVM: {svm.score(X_test, y_test):.3f}")
\end{lstlisting}

\noindent Accuracy test~: \textbf{0.968}.

% ===================================================================
\section{Stacking avec scikit-learn}
% ===================================================================

\subsection{StackingClassifier}

\begin{lstlisting}
from sklearn.ensemble import StackingClassifier

base_learners = [
    ('lr',   LogisticRegression(max_iter=5000)),
    ('knn',  KNeighborsClassifier(n_neighbors=5)),
    ('tree', DecisionTreeClassifier(max_depth=5, random_state=0)),
    ('svm',  SVC(probability=True, random_state=0))
]

meta = LogisticRegression(max_iter=5000)

stack = StackingClassifier(
    estimators=base_learners,
    final_estimator=meta,
    cv=5,
    stack_method='predict_proba'
)
stack.fit(X_train, y_train)
print(f"Stacking: {stack.score(X_test, y_test):.3f}")
\end{lstlisting}

\noindent Accuracy test~: \textbf{0.978}.

\subsection{Interpr\'etation}

Le stacking am\'eliore la performance par rapport aux classifieurs
individuels~:

\begin{center}
\begin{tabular}{lcc}
\toprule
Classifieur & Accuracy test & \'Ecart au stacking \\
\midrule
R\'egression logistique & 0.678 & $-$0.300 \\
KNN ($k=5$) & 0.918 & $-$0.060 \\
Arbre de d\'ecision & 0.942 & $-$0.036 \\
SVM (RBF) & 0.968 & $-$0.010 \\
\midrule
\textbf{Stacking} & \textbf{0.978} & --- \\
\bottomrule
\end{tabular}
\end{center}

\begin{important}
Le stacking (\textbf{0.978}) surpasse le meilleur classifieur
individuel (\textbf{SVM 0.968}). Le m\'eta-mod\`ele apprend \`a
compenser les erreurs r\'esiduelles du SVM gr\^ace aux autres
classifieurs. L'am\'elioration est modeste ici car le SVM est
d\'ej\`a tr\`es performant sur ces donn\'ees.
\end{important}

% ===================================================================
\section{Stacking manuel (compréhension)}
% ===================================================================

\`A titre p\'edagogique, on peut d\'ecomposer le stacking~:

\begin{lstlisting}
from sklearn.model_selection import KFold
import numpy as np

def stacking_manual(X, y, base_learners, meta_learner, n_folds=5):
    kf = KFold(n_splits=n_folds, shuffle=True, random_state=0)
    Z = np.zeros((len(y), len(base_learners)))
    
    # Etape 1 : predictions hors-pli (CV)
    for train_idx, val_idx in kf.split(X):
        X_fold, X_val = X[train_idx], X[val_idx]
        y_fold = y[train_idx]
        
        for k, (_, model) in enumerate(base_learners):
            model.fit(X_fold, y_fold)
            Z[val_idx, k] = model.predict_proba(X_val)[:, 1]
    
    # Etape 2 : meta-modele sur Z
    meta_learner.fit(Z, y)
    
    # Etape 3 : re-entrainer les base learners sur tout X
    for _, model in base_learners:
        model.fit(X, y)
    
    return meta_learner, base_learners
\end{lstlisting}

\noindent Cette impl\'ementation manuelle reproduit exactement le
comportement de \texttt{StackingClassifier}.

% ===================================================================
\section{Comparaison R vs Python}
% ===================================================================

\begin{center}
\begin{tabular}{lcc}
\toprule
Aspect & R & Python \\
\midrule
Stacking & \texttt{stacking::stacking()} & \texttt{ensemble.StackingClassifier} \\
M\'eta-mod\`ele & \texttt{glm} / \texttt{rpart} & \texttt{LogisticRegression} / \texttt{Ridge} \\
CV & param. \texttt{cv} (d\'efaut 5) & param. \texttt{cv} (d\'efaut 5) \\
Syntaxe & \texttt{\%\textgreater{}\%} & API orient\'ee objet \\
\bottomrule
\end{tabular}
\end{center}

\noindent En R, le package \texttt{stacking} de Tanagra offre une
interface similaire. La logique m\'etier est identique~: g\'en\'erer
les pr\'edictions niveau~1 par CV, puis appliquer le m\'eta-mod\`ele.

% ===================================================================
\newpage
\section*{Questions cl\'es (QCM)}
\addcontentsline{toc}{section}{Questions cl\'es (QCM)}

\begin{enumerate}
  \item Quelle classe scikit-learn impl\'emente le stacking~?
    \textbf{R. \texttt{StackingClassifier}.}
  \item Combien de classifieurs de base sont utilis\'es dans ce
    tutoriel~? \textbf{R. 4 (LR, KNN, Tree, SVM).}
  \item Quelle est l'accuracy du stacking en test~?
    \textbf{R. 0.978 (97.8\,\%).}
  \item Quel classifieur individuel a la meilleure performance~?
    \textbf{R. SVM (0.968).}
  \item Quelle m\'ethode g\'en\`ere les donn\'ees d'entr\'ee du
    m\'eta-mod\`ele~? \textbf{R. La validation crois\'ee (CV).}
  \item Quel m\'eta-mod\`ele est utilis\'e par d\'efaut~?
    \textbf{R. R\'egression logistique.}
  \item \`A quoi sert le param\`etre \texttt{stack\_method}~?
    \textbf{R. \`A choisir entre probas et classes pr\'edites.}
  \item Le stacking r\'eduit-il le biais ou la variance~?
    \textbf{R. Les deux (biais + variance).}
  \item Quelle est la principale diff\'erence entre Voting et
    Stacking~? \textbf{R. Stacking apprend \`a combiner.}
  \item Que fait \texttt{passthrough=True}~?
    \textbf{R. Ajoute $X$ original aux pr\'edictions niveau~1.}
\end{enumerate}

\vfill

\end{document}
